\chapter{Conclusion and Future Work}
\label{ch:sum}

\section{Conclusion}
\label{sec:conclusion}

Cross Guard was developed to solve a practical problem that web developers face regularly: checking whether the features used in their code are supported across different browsers. This process is traditionally done manually by looking up individual features on the Can I Use website, which works for small checks but becomes slow and unreliable when applied to real projects with many files and dozens of features. Cross Guard automates this entire workflow by parsing source code, detecting the web platform features it uses, and producing a detailed compatibility report.

The tool supports HTML, CSS, and JavaScript through dedicated parsers that handle the structural characteristics of each language. Together, these parsers detect hundreds of web platform features and map them to their corresponding entries in the Can I Use database. Each detected feature is checked against user-defined target browsers, and the results are used to calculate a weighted compatibility score from 0 to 100 along with a letter grade from A to F. The report breaks down results by browser, highlights unsupported and partially supported features, assigns severity levels, and provides actionable recommendations including polyfill suggestions and optional AI-generated code fixes.

The system provides two separate interfaces that share the same backend through the \texttt{AnalyzerService} facade. The desktop GUI offers an interactive workflow with drag-and-drop file upload, a visual results dashboard, analysis history, a custom rules editor, and report export. The command line interface is designed for automation and supports quality gates, multiple machine-readable output formats, and configuration generation for CI/CD pipelines. Because both interfaces use the same analysis engine, the results are always identical regardless of how the tool is used.

All analysis is performed locally, with no data sent over the network during the analysis process. The Can I Use database is bundled with the installation and can be updated on demand. The project includes a comprehensive test suite of 112 automated tests covering all major modules.

Overall, Cross Guard achieves what it set out to do. It takes a process that was manual, repetitive, and easy to forget, and turns it into something that can be done in seconds with a single command or a few clicks. It gives developers a clear picture of how compatible their code is, highlights the specific features that need attention, and suggests ways to fix them.

\section{Future Work}
\label{sec:future}

While Cross Guard in its current form provides a complete and functional tool for browser compatibility analysis, there are several areas where it could be extended and improved in the future.

\textbf{Project-Level Analysis.} Currently, Cross Guard analyzes individual files or small groups of files. A natural next step would be to support full project scanning, where the tool recursively processes all HTML, CSS, and JavaScript files in a directory. This would include respecting ignore patterns similar to \texttt{.gitignore}, aggregating results across files into a single project-level report, and identifying which features are used most frequently across the codebase.

\textbf{Framework and Preprocessor Support.} Many modern web projects do not write plain HTML, CSS, or JavaScript. Instead, they use frameworks like React, Vue, or Angular, and preprocessors like SCSS, Less, or TypeScript. Adding support for these technologies would make Cross Guard useful in a wider range of real-world projects. For example, detecting JSX syntax in React components or parsing SCSS files before analyzing the generated CSS would significantly expand the tool's coverage.

\textbf{Inline Analysis and IDE Integration.} Developers would benefit from seeing compatibility warnings directly in their code editor as they type, similar to how ESLint shows linting errors inline. Building a Language Server Protocol (LSP) integration or extensions for editors like Visual Studio Code and JetBrains IDEs would bring compatibility feedback into the development workflow without requiring the developer to run a separate tool.

\textbf{Historical Trend Analysis.} The tool already stores past analyses in a local database, but it does not currently visualize trends over time. Adding charts or graphs that show how a project's compatibility score changes across multiple analyses would help teams track their progress and identify regressions. This could also include notifications when a new analysis shows a significant drop compared to previous results.

\textbf{Browser Usage Data Integration.} Cross Guard currently weights all browsers equally by default. Integrating real-world browser usage statistics, either from global data or from the project's own analytics, would allow the scoring system to prioritize the browsers that the project's actual users rely on. A feature that is unsupported only in a browser used by less than one percent of the audience is less critical than one that fails in the most popular browser.

\textbf{Expanded Polyfill Intelligence.} The current polyfill system uses a static internal mapping between features and known polyfills. This could be enhanced by connecting to package registries like npm to check for available polyfills dynamically, estimating the size impact of adding each polyfill, and generating a ready-to-use polyfill bundle file that includes all necessary polyfills for the detected issues.

\textbf{Team Collaboration Features.} For teams working on shared projects, it would be useful to support shared configuration files that define the team's target browsers and minimum compatibility thresholds. This could be combined with a centralized dashboard where team members can view the latest analysis results for the entire project, similar to how SonarQube provides a shared code quality overview.

\textbf{Web-Based Interface.} In addition to the desktop GUI and CLI, a lightweight web interface could make the tool accessible without any local installation. Users could upload files through a browser, run the analysis on a server, and view the results online. This would lower the barrier to entry for developers who want to try the tool without installing Python or any dependencies.

The layered design with a shared backend and separate frontends makes it straightforward to add these extensions without changing the core architecture. The most impactful additions in the short term would likely be project-level scanning and IDE integration, as these address the most common workflows that developers use daily.
